\chapter[Knowledge and Empire]{How Empires Turn Knowledge into a Tool of Rule}
\section{A Suspiciously Straight Border}
Pick up an atlas.

Not the map on your phone: a paper atlas, open at Africa. Trace the Egyptian--Libyan border with a finger. It scarcely bends: a straight line running south from the Mediterranean through over a thousand kilometres of desert, as if drawn with a ruler. Look west, towards Niger, Chad and Nigeria, and many borders contain long straight segments or angular turns.

No wall or fence stands in the desert, not even a tree. The line exists on paper, and in the lives of everyone beside it. Herders born east and west may tend the same camels, speak the same language and acknowledge the same ancestors. Yet they carry different passports, pay different taxes, perform different military service and become different countries' people in war.

European powers drew these lines at negotiating tables in the late nineteenth and early twentieth centuries, using rulers and coordinates. Most draughtsmen had never visited those deserts. None of the thousands living across the lines was invited. From 1884 to 1885, European countries met for months in Berlin to discuss African spheres of influence. Not one African attended.

Notice this knowledge form's peculiarity. A map seems utterly innocent, neither shouting slogans nor wielding weapons, simply recording the world. Yet before that border became real, it was a pen stroke: drawn first, occupied and patrolled afterwards, then written into law and memorised by schoolchildren. Maps do not merely describe an existing world. They often announce a world not yet there, but about to be made.

We shall investigate these quiet announcements. Why do empires eagerly measure, classify, count, collect and name, activities we group as knowledge-seeking? What relates knowledge to power?

\section{A Pilgrim's Prayer Beads}
The time is the 1860s; the place a Himalayan pass over five thousand metres high.

Every breath feels insufficient in the thin air. Wind drives snow grains into faces like sand. A small caravan crosses the pass, led by a monk in reddish-brown robes. One hand fingers beads, the other turns a prayer wheel, while he recites. He walks slowly but remarkably steadily.

A closer observer would notice oddities.

First, the beads are wrong. A standard Tibetan rosary has 108; his has only 100.

Second, his walking is wrong. He advances one bead with each step, apparently ritually, but its rhythm follows his feet exactly, not his prayers.

Third, the prayer wheel is heavy and sounds muffled.

The ``monk'' is Nain Singh, a school headmaster from Kumaon in British India. His disguise carries several lives, including his own. Tibet then excluded foreigners; exposure could be disastrous.

The missing eight beads simplify calculation: a hundred beads count a hundred steps. Special training makes his stride almost uniform, with conversions for uphill and downhill walking. The heavy prayer wheel holds not scripture but slips recording dates, routes, paces and directions. At night, when companions sleep, he writes observations secretly gathered during the day. Stars supply latitude; boiling water supplies altitude. At elevation, water boils below one hundred degrees Celsius; a lower boiling point means greater height. This was an accepted surveying method.

For whom does he work? The British. Beginning in the early nineteenth century, British India undertook the enormous Great Trigonometrical Survey. Hundreds of workers hauled theodolites weighing tens of tonnes across mountains, mapping the subcontinent precisely mile by mile. Lasting decades, it was among humanity's largest contemporary scientific projects. But it stalled at the Himalayas: instruments crowded one side; the other was inaccessible. Survey officers recruited Indians, training them as disguised pilgrims and merchants to walk maps into existence through bodies and memories. These men were called pundits, meaning scholars. Nain Singh was the best known: he reached Lhasa and returned alive.

Stay at this pass a little longer and examine its contradictions.

A colonised man gathers intelligence for colonisers using precisely the knowledge his local identity provides: languages, appearance, religious practices, passable roads and villages offering lodging. British officers possess excellent instruments, none replacing his face or feet. When maps appear, however, the survey department and officers receive the bylines; men like Nain Singh hide in footnotes.

What is he? An imperial tool, or an explorer accomplishing his own feat with imperial money? An exploited knowledge source, or its real author? Wait before answering. We shall return to the pass.

\section{How Can Knowledge Be Guilty?}
Descend and set out the conflict.

On one side stand impeccably reasonable activities. Maps locate roads and flood-prone rivers. Population counts guide taxes and grain reserves. Plant collecting investigates medicines. Dictionaries and grammars preserve disappearing languages. Museums preserve ancestral craftsmanship. Who opposes knowledge? Would that not mean supporting ignorance?

On the other stand disturbing facts. The Great Trigonometrical Survey advanced alongside Britain's conquest of India: armies and tax officers followed its maps. Late-nineteenth-century censuses asked households their caste and printed answers in official tables. Previously ambiguous, mobile identities varying by locality hardened into lifelong or hereditary labels. In Belgian-ruled Rwanda, colonial authorities registered ethnic identities for administration, issuing cards marked Hutu or Tutsi. Decades later, during the 1994 genocide, killers inspected such cards at checkpoints; one line decided life or death. Consider plants too. Andean residents already knew cinchona bark treated malaria. In the nineteenth century, Europeans exported seeds to Indian and South-East Asian plantations. Extracted quinine first enabled long-term tropical military occupation: life-saving medicine also equipped imperial expansion. Colonial language surveys recorded hundreds of languages and dialects, still a linguistic treasure. Their immediate uses included identifying who spoke what, drawing administrative districts, choosing official languages and training interpreters and police.

The real question surfaces:

\textbf{Is knowledge neutral?}

Are maps, tables, specimens and racial classifications faithful records or shapes cut to particular needs? If the latter, should we call such knowledge false? Wait: Nain Singh's routes were accurate; quinine really reduced fever; the survey's Everest height closely resembled modern satellite measurements. Accuracy and troubling implications can coexist.

That coexistence is the gap we shall enter.

\section{What the Viceroy and Villager Saw}
The same census form becomes different things on a Delhi viceregal desk and an Indian village threshing floor. Present both positions fully.

\textbf{Position one: a public undertaking.}

Make the imperial surveyor's, statistician's and botanist's case seriously. Once ``knowledge serves rule'' becomes automatic, it can slide into calling all maps and statistics conspiracies, ending analysis.

Their argument has at least four levels. First, accurate public information saves lives. Maps enable irrigation, railways and relief; statistics provide standards for aid and taxes. What standards replaced was no pastoral idyll, but arbitrary tax demands and guessed famine deaths. Second, medicine's account extends beyond soldiers. Quinine served colonial armies, but more recipients were ordinary tropical residents, including colonised patients. Antimalarials still save many lives annually. Third, genuine curiosity existed. Linguists and botanists sincerely admired studied cultures, risking lives to record languages and species. Their archives now sometimes provide endangered cultures' only surviving documentation. Fourth, motives cannot cancel consequences. Whatever their original purposes, maps, statistics and medicine remained as infrastructure for independent successor states.

These are strong reasons. To challenge them, acknowledge that you dispute something other than knowledge's usefulness.

\textbf{Position two: turning living people into files.}

From the threshing floor, matters look different.

``What is your caste?'' sounds neutral. Yet in many places answers once depended on context: farming, marriage and urban life could involve different identities. Identity flowed like water, shaped by its container. Forms accept boxes, not water. Officials compelled a choice or filled one themselves. Generations later, boxes became real. People understood themselves through printed identities, organised parties, applied for quotas and divided allies from enemies accordingly. Classification ceased merely describing populations; it produced them.

Maps did likewise. Woodland around a village might be shared for grazing, firewood and worship, without clearly articulated boundaries because such boundaries did not exist. Surveyors arrived with stakes, lines and registers, turning commons into state forest or private estates. Yesterday's sheep path became today's trespass. Knowledge did not lie; it inscribed one institution's answers into the ground in advance.

Naming mattered too. Nepalese called the highest mountain Sagarmatha, Tibetans Chomolungma, each with ancient meanings. Yet for over a century, world maps named it Everest, after a British survey director. Botanical gardens replaced plants' centuries-old local names with Latin names, crediting European collectors as discoverers instead of the guides teaching their uses. Museums were more direct: a ceremonial sculpture crossed the sea into a glass case labelled with its donor and ``nineteenth-century West African art''. For its maker, it was not art but an ancestor.

\textbf{A third voice: do not reduce people to passive information sources.}

The account above contains a trap: empire collects, Indigenous people are collected. Empire is subject, local people object. Our Himalayan scene warns against that simplicity.

Colonial knowledge machinery depended throughout on guides, interpreters, informants, collectors, clerks and surveyors like Nain Singh. They were no taps to be turned for information. They selected, concealed, misled and bargained. Sacred sites excluded from maps, roads avoided, words translated or withheld: these were real decisions. Confident colonial reports of local customs often trace back to whatever an informant wanted an official to believe. Empire imagined itself painting colonial portraits; colonised people held at least half the pen.

More striking was the counterstroke. Dadabhai Naoroji collected British India's own trade and financial statistics, calculating how Indian wealth systematically flowed to Britain. His collected analysis appeared in 1901, known as the drain theory. Notice the act: he did not destroy the statistical machine. He took its numbers, turned its barrel and fired at its builders. Indian nationalists repeatedly did likewise, using colonial censuses to claim parliamentary seats and railway maps to organise nationwide mobilisation. Knowledge can serve power, but also be captured, modified and used against it. We shall develop this point.

\section[Thinking Bridge: Knowledge's Origins]{Thinking Bridge: Run a Background Check on Knowledge}
Can we judge maps, tables, classifications, specimens and statistics through a portable tool rather than intuition? Yes: investigate any piece of knowledge's background with four questions.

\textbf{First: who made it?}

Not merely the byline: who paid, chose the question and received the result? The East India Company and later colonial government funded the Great Trigonometrical Survey; armies and revenue offices received it first. That origin does not invalidate maps, but reminds us they were not initially made for village shepherds.

\textbf{Second: which categories cut it up?}

Every count divides a continuous world into boxes. Other divisions are always possible: people by occupation rather than caste, land by use rather than ownership. Who chose? What was removed, merged or ignored?

\textbf{Third: who was absent?}

Could measured people speak, inspect, challenge or alter their records? No Africans sat at Berlin's table; census forms had no ``I cannot say what I count as'' option; specimen labels omitted guides. Absence itself supplies structural information.

\textbf{Fourth: what did it make possible?}

Most crucially, which impossibilities became possibilities? Registers made taxation and conscription precise; ethnic identity cards made checkpoint screening precise; ethnic maps sharpened divide-and-rule. Follow this question to see knowledge's conduit into power.

In Seeing Like a State, 1998, political anthropologist James C. Scott offered a vivid formulation. States deliberately simplify complex, fluid realities into orderly, readable boxes for administrative convenience, pursuing what he called legibility. Remember his qualification: boxes are always simpler than reality. Governing the boxes as though they were reality often brings disaster. The map is not the territory; the table is not the population. Put that beside every statistic and you have the tool.

Practise nearby. A class ranking: who made it, why classify only by marks, whose opinions are missing, and what does it enable---class allocation, awards, a parent-meeting conversation? Ask the same four questions of a nutrition label, an app's personality report or a shop's customer profile. Many things will reveal their shape.

\section{Maps and Registers Three Millennia Ago}
Read a primary source. Linking knowledge and rule predates modern Europe by a great deal.

The Rites of Zhou, compiled between the Warring States and Han periods, idealises and largely designs Zhou government. Its section on the Minister of the Earth opens the Grand Minister's responsibilities thus:

\begin{quote}
The Grand Minister's duty is to administer the maps of the state's lands and the numbers of its people, assisting the king in settling and governing the realms.
\end{quote}
---Rites of Zhou, ``Minister of the Earth''.

In other words, the official holds land maps and population counts to help the ruler govern securely.

Unpack the sentence. More than two millennia ago, this ideal bureaucratic job description placed \textbf{maps} and \textbf{population figures} first. Only afterwards could governing proceed. Sequence itself makes the argument: knowledge precedes rule and provides its foundation.

This intuition was neither uniquely Chinese nor merely ancient. The Book of Documents' Tribute of Yu divides the world into nine provinces, listing soils, products and tribute, resembling an early nationwide resource survey. Scholars generally date it to the Warring States, summarising contemporary geography under Yu's name; we treat that as an attributed account, not Yu's handwriting. Ming Yellow Registers recorded households; fish-scale cadastral maps registered fields, each scale-like parcel labelled with owner and area. Taxes, lawsuits and conscription relied on them. Across Eurasia, Napoleon took over a hundred scholars on his Egyptian expedition in 1798. They surveyed, documented antiquities and collected flora and fauna, later compiled in the massive Description of Egypt. Occupation forces and scholars shared ships: conquest and knowledge production held the same ticket, with an openness almost honest in itself.

Placing Zhou bureaucracy beside Napoleon gives an important reminder: \textbf{knowledge as a ruling tool is not one civilisation's peculiar vice, but a near-universal practice of large states}. Chinese emperors and European viceroys were strikingly alike. Our question is not why Westerners were bad, but something general: once government exceeds a society of personal acquaintance, it must turn distant lives into paper numbers and diagrams. That conversion is never merely technical.

\section{Why Empires Measure: Three Accounts}
We have enough components for genuinely different explanations. First distinguish four categories. Empires mapped land, counted people, collected objects and compiled dictionaries: evidence. Why they did so invites competing accounts. Officials' belief in civilisation and progress records their own understanding. Our evaluation is a value judgement. Here we compare explanations.

\textbf{Explanation one: administrative technology---a large machine needs a dashboard.}

Govern three people from memory; govern three hundred million through archives. Maps, statistics and registers are imperial dashboards and drive shafts, like factory inventories. This concise explanation addresses a cross-cultural fact: Chinese, Ottoman and British empires all did it, as do ideologically different states. It is a consequence of scale, not one culture's conspiracy. But why is the dashboard skewed? If knowledge alone is the aim, why classify by caste and ethnicity rather than occupation? Why ship specimens to royal gardens instead of retaining them locally? Why advanced instruments but absent guides' names? Administrative need cannot explain the consistent direction: knowledge flows to capitals, conclusions favour taxation and conscription.

\textbf{Explanation two: power and knowledge are one---knowledge itself produces power.}

Michel Foucault proposed that these are not separate things: power does not simply exploit innocent knowledge. It determines worthwhile questions, acceptable evidence and recognised expertise; knowledge in turn produces the social reality power seeks. India's census did not merely record pre-existing castes. It cast fluid identities into fixed compartments, around which people then organised and opposed one another. \textbf{Knowledge described a reality it helped manufacture.} This sharp account illuminates naming, classification's effects and absent guides. Its danger is excessive sharpness: if all knowledge is entangled with power, how distinguish an accurate map from propaganda? Why call Everest's measured height true and innate ethnic inferiority false? Treat everything as power's product and criticism's blade becomes blunt.

\textbf{Explanation three: curiosity and honour---mixed motives, consequential outcomes.}

Do not imagine uniform people. Some surveyors loved geometry, some botanists thrilled at a leaf, some missionary lexicographers sincerely sought language preservation, and Nain Singh had his own adventurousness and pride. Empire worked precisely by recruiting not identical cold bureaucrats but individuals with varied passions. This matches lived human complexity and explains surveying precision: only people truly devoted to measurement could achieve it. But motives do not explain direction. Why did individual enthusiasms converge on imperial usefulness? Who controlled funding, permissions, publishing and rewards? Pure curiosity was real; so were the gates.

Together these accounts make useful spectacles. Scale compelled knowledge-seeking; power shaped categories and flows; individuals supplied diverse motives within the machine.

\section{Further Challenge: Who Represents Whom?}
Our deepest theoretical question comes from postcolonial studies.

In Orientalism, 1978, Palestinian American scholar Edward Said examined centuries of European scholarship, literature, travel writing and official reports about the East, including the Middle East and Asia. Broadly, he found a shared frame beneath learning and goodwill: an East depicted as mysterious, static, irrational and awaiting understanding, opposite a rational, progressive Europe entitled to understand it. The key was not whether one book was accurate, but that \textbf{a society's long monopoly on describing another makes description itself part of domination}. Who studies, who becomes specimen? Who speaks, who is spoken about? This asymmetry runs deeper than individual prejudice.

Said pinned a question to the humanities' wall: \textbf{who represents whom?} Notice representation's double meaning: depicting, as a portrait represents someone, and speaking for, as an MP represents someone. Postcolonial insight joins the two. A museum label describes a sculpture while speaking for its makers, whose descendants cannot review it. A textbook calling a people fatalistic both describes and stamps a judgement on millions. Every representation implies authorisation: who authorised you?

Once asked, the question cannot be withdrawn. It challenges textbooks representing history, documentaries representing distant places, news representing participants and experts representing publics. It even bites postcolonial studies, as an honest theory must allow. In 1988, Gayatri Spivak published ``Can the Subaltern Speak?'' When Western or local intellectuals righteously speak for marginalised people, do they drown out those people's voices? Can representation itself dispossess? Her answer was pessimistic, and debate continues. Here we relay her question without deciding it for her, perhaps a small act of respect.

Return to the pass. Nain Singh walked Lhasa's routes, but officers' names appeared on maps. A real knowledge producer was registered as a local assistant. Representation's asymmetry is plain. Go further without reducing him to victimhood. He deceived imperial adversaries with his disguise and partly used imperial resources, receiving pay, recognition and a remarkable life. Colonised people were not passive repositories awaiting harvest. They negotiated, collaborated, resisted, hedged and co-authored knowledge. Acknowledging agency does not excuse imperial control over attribution. It restores people as calculating, choosing, contradictory humans. A history of pure victims makes the same mistake as colonial specimen archives: neither sees full persons.

Naoroji's redirected statistical weapon gains depth here too. Postcolonial study does not demand destroying knowledge, but asking who holds it and for whom it speaks. Knowledge can be captured. Maps can be redrawn, museums reorganised and dictionaries continued by speakers themselves. Oppressed people's strongest reclaimed weapon is often not a blade but the right to define themselves.

\section[The Survey in Your Phone]{Today: The ``Great Trigonometrical Survey'' in Your Phone}
This old craft of knowledge-based government has not entered a museum. It exceeds empires' wildest dreams in scale, while replacing land with you as its object.

Every tap, search or extra two seconds on a video walks another step for an invisible surveying machine. Platforms' user profiles are census forms: age, gender, location, purchasing power, political tastes and emotional tendencies divided into thousands of boxes. The logic resembles India's census: fix your fluid self into a classified self, then target advertisements, videos and prices accordingly. Measured people are absent, unable to inspect, much less change, the record. Our four questions remain powerful: who made it, how did they divide it, who cannot speak, and what does it enable?

Consumer genetic testing brings racial classification closer. Companies report percentages of Northern European or East Asian ancestry. Scientific-sounding figures compare your DNA segments with company-defined reference populations; different companies draw different groups and produce different answers. Old racial classification measured skulls; new systems compare genes algorithmically. The tools differ enormously, but an institution still defines population boxes and assigns you. The boxes are real in that calculations occurred; their boundaries remain human choices.

Museum debts are being recalculated too. In 1897, British forces conquered the Kingdom of Benin in present-day Nigeria, looting thousands of palace bronzes and ivories that dispersed through Western museums as the Benin Bronzes. Since independence, Nigeria has pursued them for over a century. Recently, some German, British and American institutions have begun returns. Greece continues seeking the British Museum's Parthenon sculptures. Opposing arguments miniaturise this chapter: ``We preserve them better for all humanity'' versus ``These are stolen ancestors and dignity, not exhibits; who authorised your representation of humanity?'' Said's question echoes through galleries.

Resistance upgrades too. Today's Naorojis use platforms' data against them: researchers scrape recommendation patterns, public-interest groups use satellite maps to track deforestation on Indigenous land, and minority communities write their histories and rescue languages through social media. Empires once compiled dictionaries; communities now compile their own. Knowledge and power have no final settlement, only successive rounds.

\section[Your Turn: Knowledge without Power?]{Your Turn: Can Knowledge Be Untouched by Power?}
Return to the atlas.

You now know that straight lines were drawn, not discovered; census boxes were cut, not found; museum labels represent makers who were absent, not neutral descriptions.

You also know the other half: Nain Singh measured accurately, quinine reduces fever, archives remain linguistic resources, and people entered in tables may seize them to calculate imperial accounts and write their histories.

Answer our final question in your own way, with reasons, not a standard formula:

\textbf{If you someday hold a knowledge machine---designing a questionnaire, running recommendations, managing a museum or writing a textbook---can you keep it free of power?}

One answer says no. Classification, selection and speaking for others already involve power. Pretended neutrality merely conceals it. Honesty therefore requires declared positions and open scrutiny: let measured people inspect, challenge and amend records, and represented people reclaim the microphone.

Another asks whether this makes all knowledge guilty, requiring an end to counting, mapping and museums. But epidemics require statistics, relief requires maps, endangered languages require records. Abandoning inquiry because power touches it may abandon knowledge alone. Power remains; nobody produces knowledge.

Which do you choose? Or is the question wrong? Perhaps ask not whether knowledge touches power, but how distant measurers and measured people stand before the machine, and whether that distance can shrink.

Before answering, look once more towards the pass. The man fingering false prayer beads is long gone; snow holds no footprints. Yet in every map, form and statement made for you, someone still walks his hundred steps, deciding which numbers to surrender and which to retain inside a prayer wheel. You are the next person measured. You are also the next measurer.
